\cotefile{161-2}
\scanpage{96}
La catégorie (ou \unsure{un} \uline{lieu} quelc.) universelle
pour ceci est la catégorie des faisceaux sur
$R'$ pour la top. de Zariski (= foncteurs
contravariants $R\longrightarrow\mathrm{Ens}$ qui transforment
sommes en produits). Il se trouve que
c'est un topos \deleted{$\unsure{E|\mathcal E|\circledast}$} Le corps universel est
le faisceau structural évident sur ce topos....

\centerline{\rule{5em}{0.4pt}}

$\overline{\mathrm M}$ question avec \uline{anneau intègre}. Soit
$O'\subset O$ le sous-objet ``qui opère régulièrement sur $O$'',
$(O'\supset O^*)$ est il définissable par $\varprojlim$ et $\varinjlim$ ? Commute-t'il aux images inverses \unsure{par un morphisme} de topos ?

\deleted{$O'$ définissable par $\varprojlim$ \unsure{finies} ? Si $\unsure{E}=\unsure{\mathrm{Sch.aff}|\mathrm{Sch.ff}}^{\unsure{\sim}}$, cela
signifie que les foncteurs \unsure{fonctionnant} sur Sch aff
Cov. sur \unsure{Ann} qui : \unsure{tt}}

Dans le cas \unsure{de|où|des} l'anneau universel d'un
topos, $O\in\widehat R$, $O'(\deleted{\unsure{A|X}})$ est le sous-\deleted{\unsure{anneau}}\inserted{ens de}
de $A=O(X)$ formé des $f\in A$ tels que $f$ soit
régulier, et \uline{le reste} par tt chgt d'anneau
de type fini $A\longrightarrow A'$ -- mais cela \deleted{signifie}
implique que $f\notin m$ pour $m\in\mathrm{Max}\,A$, donc
$f$ inversible. Donc \unsure{on a} $O'=O^*$ ! Or
un foncteur fibre sur $\widehat R$ est un anneau
(ensembliste) et il n'est pas vrai qu'un tel
anneau soit intègre \unsure{ssi} c'est un corps, donc
on n'a pas $O'_p=O_p^{*}$ i.e $A'=A^*$ !

% archive sheet number: 95
